\begin{enumerate}
    \item tableau d'avancement~:

    \begin{tabular}{|l|l|l|l|l|l|}
    \hline \multicolumn{2}{|c|}{Equation de la réaction} & \multicolumn{4}{|c|}{\begin{tabular}{l}
    $$
    \mathrm{C}_6 \mathrm{H}_5 \mathrm{COOH}+\mathrm{H}_2 \mathrm{O} \longrightarrow \mathrm{O}_6 \mathrm{H}_{(i)} \mathrm{COO}^{-}+{\underset{(\text { al })}{\mathrm{H}_3 \mathrm{O}^{+}}}^{+}
    $$
     \\
    \item \\
    \item \\
    \item \end{tabular}} \\
    \hline Les états & avancement & $n\left(\mathrm{CH}_3 \mathrm{COOH}\right)$ & $n\left(H_2 \mathrm{O}\right)$ & $n\left(\mathrm{CH}_3 \mathrm{COO}^{-}\right)$ & $n\left(\mathrm{H}_3 \mathrm{O}^{+}\right)$ \\
    \hline état inttal & 0 & $\mathrm{n}_0$ & excès & 0 & 0 \\
    \hline état de transformation & $x$ & $\mathrm{n}_0-x$ & excés & $x$ & $x$ \\
    \hline état final & $x_f$ & $\mathrm{n}_0-x_f$ & excés & $x_f$ & $x_f$ \\
    \hline
    \end{tabular}
    \item Or la conductivité se mesure lorsque l'état final est atteint :

    $$
    \sigma=\lambda_{\left(\mathrm{CH}_1, \mathrm{COO}^{-}\right)^{-}}\left[\mathrm{C}_6 \mathrm{H}_5 \mathrm{COO}^{-}\right]_f+\lambda_{\left(\mathrm{HO}^{+}\right)^{-}}\left[\mathrm{H}_3 \mathrm{O}^{+}\right]_f
    $$

    D'après le tableau d'avancement on a:

    $$
    { }_f^n\left(\mathrm{H}_3 \mathrm{O}^{+}\right)=n\left(\mathrm{C}_6 \mathrm{H}_5 \mathrm{COO}^{-}\right)=x_f
    $$


    $$
    \text { done: } \quad\left[\mathrm{H}_3 \mathrm{O}^{+}\right]=\left[\mathrm{C}_6 \mathrm{H}_5 \mathrm{COO}^{-}\right]=\frac{x_f}{V} \quad \Rightarrow \quad \sigma=\left(\lambda_{\mathrm{T} B} \mathrm{COO} O^{-}+\lambda_{\mathrm{HO}}+\right) \cdot \frac{x_f}{V}
    $$
    \item l'avancement final de la dissociation de l'acide benzoïque dans l'eau:

    $$
    \begin{gathered}
    x_f=\frac{\sigma \cdot V}{\lambda_{\left(c_f f_f, 000^{-}\right)}+\lambda_{\left(H_3 O^{+}\right)}} \\
    x_f=\frac{36,1 \cdot 10^{-3} S m^{-1} \cdot 50 \cdot 10^{-6} m^3}{(35+3,23) \cdot 10^{-3} S \cdot m^2 m o l^{-1}}=4,72 \cdot 10^{-5} \mathrm{~mol}
    \end{gathered}
    $$
    \item $\left[{ }^{\circ} \mathrm{C}_6 \mathrm{H}_5 \mathrm{COO}^{-}\right]_f=\left[\mathrm{H}_3 \mathrm{O}^*\right]_f=\frac{x_f}{V}=\frac{4.72 \cdot 10^{-5} \mathrm{~mol}}{0,05 \mathrm{~L}}=0,94 \cdot 10^{-3} \mathrm{~mol} / \mathrm{L}$
    \item $p H=-\log \left[H_3 O^{+}\right]=-\log \left(0,94 \cdot 10^{-3}\right)=3 \quad$ car le pH se mesure lorsque l'état final est atteint
    \item $\tau=\frac{x_f}{x_{\max }}=\frac{x_f}{C V}=\frac{4,72 \times 10^{-5}}{1,18 \times 10^{-2} \times 50 \times 10^{-3}}=0,08=8 \% \quad \tau<1 \Rightarrow$ la réaction est limitée.

    Cela signifie que seulement $8 \%$ des molécules d'acide benzoïque ont été transformées pour donner leur base conjuguée et $\mathrm{H}_3 \mathrm{O}^{+}$.
\end{enumerate}

